% Pursuing Stacks (cotes 134-2 to 134-8): archive cote 134-7, site page 45 (PDF page 46), typescript page 487.
% typescript. Draft compared: ps6.tex:2215-2280.
% AI-transcribed by Claude Opus 5.5 (claude-opus-5-5, adaptive thinking, effort high), Message Batches API, batch ps3 (msgbatch_01XdFmuVVBAcgc7EFxdtCqsj), 2026-11-12; not yet checked by eye.
\scanpage{134-7/45}
\mspage{487}

for any unipotent bundle X over k,  and any quasicoherent sheaf
\[ F = W(N) \quad , \]
where N is any k-module.  Of course, (19)  yields (for general $\Gamma$ )
\[ (19') \qquad \mathrm{Ext}^{i}_{\underline{O}_{\hand{k}}}( \Gamma , F) \simeq \varinjlim_{\hand{\alpha}} \mathrm{Ext}^{i}_{\underline{O}_k}( \Gamma_{\hand{\alpha}} , F) \quad , \]
so that (20')  may be rewritten more explicitly,  if  $X \hand{\simeq} W(M)$, \hand{as}
\[ (21) \qquad H^{n}(X,W(N)_X) \hand{\simeq} \xout{\mathrm{lim}}_{\deleted{\hand{i}}} \hand{\bigoplus_{i\geq 0}} \mathrm{Ext}^{n}_{k}(\Gamma_k^{\ i}(M),N) \quad . \]
At any rate,  we have a canonical map (20) in $D^{+}(Ab_k)$,  hence maps
(20'),  in view of the isomorphism (5)  and the canonical map
\[ (\hand{*}) \qquad \underline{O}_k^{(X)} \hand{\longrightarrow} L(X) \]
deduced from (18),  and the question now is whether th\over{is}{ese} \over{is}{are} isomor-%
phism\hand{s}. \deleted{The question \inserted{(20)} makes a sense for any X,  we are more specifically
interested though in the case when X is a unipotent bundle,  to which
\xout{We may of course assume X = W(M)}
case I will restrict.} We may of course assume  in (21)
\[ X = W(M) \quad , \quad n \hand{\geq} 1 \ . \]
If M is projective,  so are the modules $\Gamma_k^{\ i}(M)$,  and hence \xout{\unsure{for n>0}}
the second  member in (21) is zero,  so we should check the first mem-%
ber is too.  This \over{\unsure{a}}{i}s clear when M is of finite type,  hence X represen-%
table by an affine scheme,  who\over{\unsure{e}}{s}e quasi\over{\unsure{o}}{c}oherent cohomology is well known
therefore to vanish in dim. $n \hand{>} 0$.  T\hand{h}e  general case should be a con-%
sequence of this,  representing X as the filtering direct limit of
\over{\unsure{p}}{i}t's submodules which are projective of finite type  -  this should work
at any rate when M is free with a basis which is at most countable,
using the standard so-called ``Mittag\deleted{\unsure{e}}-Leffler'' argument for passage to
limit.   Thus in case M projective,  (21)  and hence (20) seems OK
indeed.  When M is \uline{not} projective however,  \over{\unsure{w}}{t}here must be some k-module
N such that  $\mathrm{Ext}^{1}_{k}(M,N) \neq 0$ ,  and hence the second member of (19) is
non zero for n=1,  which should imply rather unexpectedly
\[ H^{1}(X,F) = H^{1}(W(M), W(N)_{W(M)}) \neq 0 \quad , \]
whereas till this very moment I \over{has}{had} \inserted{been} under the impression that quasi-coh\deleted{e}
rent cohomology of quasiunipotent bundles should be zero,  just as
for affine schemes !  Maybe it has been familiar to Larry Breen for a
long time that this is \uline{not} so\over{!}{?}  Maybe \inserted{also} for what we want to do it is'nt
really basic to find out whether (21) is true in full generality,  as
for the purpose of studying P\over{\unsure{a}}{o}stnikov type dévissage,  the \deleted{coefficients}

\margin{45}
